% only used acronyms are displayed
% \ac{LOL}  -> First time: Laughing Out Loud (LOL), Else: LOL
% \acs{LOL} -> LOL
% \acf{LOL} -> Laughing Out Loud (LOL)
% \acs{LOL} -> Laughing Out Loud
% \acp{LOL} -> LOLs
% documentation: https://www.ctan.org/pkg/acronym

\chapter*{Liste des abréviations}

\begin{acronym}[AABBCCD]  			% define width of left column 
\setlength{\itemsep}{-\parsep}			% define spacing between acronyms

% i can add my own acronym that i used in the docs 



% Géométrie et données de l'avion
\acro{Sw}{Surface alaire}
\acro{bw}{Envergure de l'aile}
\acro{hw}{Hauteur de l'aile au-dessus du sol}
\acro{W}{Poids de l'avion}
\acro{Ra}{Coefficient de rapport d'aspect}
\acro{lambda}{Paramètre de rapport d'aspect}

% Coefficients aérodynamiques
\acro{Cdo}{Coefficient de traînée à portance nulle}
\acro{Cdol}{Terme additionnel du coefficient de traînée}
\acro{Cd}{Coefficient de traînée}
\acro{Cl}{Coefficient de portance}
\acro{Clmax}{Coefficient de portance maximal}
\acro{e}{Facteur d'efficacité d'Oswald}

% Paramètres du mouvement au sol
\acro{mu}{Coefficient de frottement au sol}
\acro{Vlo}{Vitesse de décollage}
\acro{Vhw}{Vitesse de référence}
\acro{Vi}{Vitesse initiale pour l'intégration}
\acro{Vip1}{Vitesse finale pour l'intégration}
\acro{g}{Accélération de la pesanteur}
\acro{rho}{Masse volumique de l'air}

% Paramètres de propulsion
\acro{P}{Puissance du moteur à pistons}
\acro{T}{Poussée du moteur à réaction}
\acro{efficiency}{Rendement de l'hélice}
\acro{T0}{Poussée initiale ou coefficient expérimental de poussée}
\acro{T1}{Coefficient expérimental de poussée}
\acro{T2}{Coefficient expérimental de poussée}
\acro{A1}{Coefficient de la formule de poussée}
\acro{A2}{Coefficient de la formule de poussée}
\acro{A3}{Coefficient de la formule de poussée}
\acro{A4}{Coefficient de la formule de poussée}

% Paramètre numérique
\acro{n}{Nombre d'itérations}


\end{acronym}
